% !PRINTED 122
\weberpage{122}
%<<CONT

and which is also \gesperrt{continuous} in this interval. In agreement with \S\,31, V and VI, this means that
\begin{equation}
 f(x\pm h)-f(x)
\tag{2}
\end{equation}
falls, in absolute value, below an arbitrarily prescribed number $\eta$ throughout the interval whenever the positive $h$ is smaller than a certain number $\varepsilon$. For $x=a$ we take in (2) only the upper sign, and for $x=b$ only the lower sign, so as not to pass outside the interval with $x\pm h$.

For the values of such a function in the interval there is, by 1., a \gesperrt{lower bound} $g$, and we now prove the proposition:

3. \gesperrt{The function $f(x)$ assumes the value $g$ for some value $\xi$ of the interval, so that the lower bound becomes a minimum of the function $f(x)$ in the interval.}

The proof is as follows. By 2., the \gesperrt{lower bound} of the function values in a part of the interval (1) is either equal to or greater than $g$.

If now $f(a)=g$, what is to be proved is true. If, however, $f(a)>g$, then by the continuity of $f(x)$ one can give an interval $a\le x\le a+h$ in which $f(x)>g$ remains true, and hence the \gesperrt{lower bound} of $f(x)$ is greater than $g$.

We now construct in the interval (1) a cut $(\mathfrak A,\mathfrak B)$ of the following kind: we count a value $\alpha$ of the interval (1) as belonging to $\mathfrak A$ if the \gesperrt{lower bound} of $f(x)$ in the interval $a\le x\le \alpha$ is greater than $g$, and a value $\beta$ as belonging to $\mathfrak B$ if the \gesperrt{lower bound} in the interval $a\le x\le \beta$ is equal to $g$.

It is possible that $\mathfrak B$ consists only of the single value $b$; but then $f(b)=g$ must hold, and our assertion is fulfilled for $x=b$. For if $f(b)>g$, one could choose a quantity $g'$ between $f(b)$ and $g$, and, by continuity, an interval $b-h\le x\le b$ such that in this interval all function values $f(x)$ are greater than $g'$, so that their \gesperrt{lower bound} is equal to or greater than $g'$, and hence certainly greater than $g$. But since $b-h$ belongs to $\mathfrak A$, the \gesperrt{lower bound} is also greater than $g$ in the interval $a\le x\le b-h$, and consequently in the whole interval (1), contrary to the assumption.
